\sectionguard
\section{The Boundary of the Defined: Titles, Doctrines, Opinions}

\subsection*{Authoritative teaching that is not a fifth dogma}

Several firm Marian teachings stand near the four dogmas without being
among them, and a reference must keep the categories visible. Mary's
freedom from personal sin throughout her life is constant authoritative
teaching (CCC 493: ``By the grace of God Mary remained free of every
personal sin her whole life long'') but is not the object of the 1854
definition. Her spiritual motherhood of the faithful and the title
Mother of the Church (proclaimed by Paul VI in 1964, echoed at CCC 963,
975) are authoritative teaching. Her Queenship is taught by the
Assumption definition's own horizon and by LG 59 (``exalted by the Lord
as Queen of the universe''), yet no separate coronation dogma exists.
Each of these binds in its own mode; none carries the four dogmas'
definitional form.\endnote{CCC 493, 963--975, read at the Vatican
IntraText pages cited above; LG 59 at the Vatican English text of
\work{Lumen gentium}.}

\subsection*{Cooperation, mediation, and the disputed titles}

Mary's cooperation in the work of salvation is real and conciliar
doctrine: she ``cooperated by her obedience, faith, hope, and burning
charity in the Savior's work,'' and is ``a mother to us in the order of
grace'' (LG 61, quoted at CCC 968). But the Council fenced the language
on every side. There is ``but one Mediator'' (1 Tim 2:5; LG 60), and
Mary's maternal role ``in no wise obscures or diminishes this unique
mediation of Christ, but rather shows His power''; her influence
``flows forth from the superabundance of the merits of Christ, rests on
His mediation, depends entirely on it and draws all its power from it''
(LG 60). The Church does invoke her ``under the titles of Advocate,
Auxiliatrix, Adjutrix, and Mediatrix,'' yet ``this, however, is to be so
understood that it neither takes away from nor adds anything to the
dignity and efficaciousness of Christ the one Mediator'' (LG
62).\endnote{LG 60--62, checked at the Vatican English text; CCC
968--970 quotes the same passages.}

The proposal to define a fifth Marian dogma---Mary as Co-redemptrix,
Mediatrix of All Graces, and Advocate---has circulated in petition
movements since the twentieth century. It stands \emph{outside} the four
dogmas, and the Holy See has now judged its central title directly. The
Dicastery for the Doctrine of the Faith's doctrinal note \work{Mater
Populi Fidelis}, approved by Pope Leo XIV on 7 October 2025 and published
4 November 2025, concludes that ``it is always inappropriate to use the
title `Co-redemptrix' to define Mary's cooperation,'' because the title
``risks obscuring Christ's unique salvific mediation'' (\S22). On
mediation it teaches that the Council deliberately preferred the language
of cooperation and maternal assistance (\S27); on ``Mediatrix of All
Graces'' it observes that Mary, ``the first redeemed, could not have been
the mediatrix of the grace that she herself received'' (\S67), while
allowing that her intercession may implore the ``actual graces'' that
prepare and accompany conversion (\S69).\endnote{\work{Mater Populi
Fidelis}, \S\S22, 27, 67, 69, read at
\url{https://www.vatican.va/roman_curia/congregations/cfaith/documents/rc_ddf_doc_20251104_mater-populi-fidelis_en.html}.
Historical papal uses of ``Co-redemptrix'' (reviewed at \S18) retain
their own dates and authority; the note governs present constructive
use. The note is a papally approved doctrinal note, not a solemn
definition; it is the current controlling statement on these titles.} \term{Advocate}, retained by LG 62, means Mary's maternal
intercession within Christ's mediation---never protection from a
merciless Father or a reluctant Christ.

The doctrinal test that governs every Marian title follows from the four
dogmas themselves: whatever is said of Mary must be received, created,
participated, and wholly dependent on the one Redeemer. Said that way,
the titles guard the dogmas; said otherwise, they obscure them. That is
the boundary this reference draws, and it is the Church's own.
